\section{Proofs of Propositions~\ref{prop:mean}--\ref{prop:median}}
\label{app:proofs}

Throughout, $b=\tau+\varepsilon$ and $H_t=c_t-x_t$.

\paragraph{Proposition~\ref{prop:mean} (mean).} The mean is linear: $\bar y\ge b$ iff $\sum_t(y_t-x_t)\ge nb-\sum_tx_t$, and any $y\in[x,c]$ with that total raise is feasible, which requires $\sum_tH_t\ge nb-\sum_tx_t$ and gives $L_1^\ast=\max(0,nb-\sum_tx_t)$. $K$ is minimised by filling the largest headrooms first. If no cap binds, a uniform raise $d$ gives mean $\bar x+d$, so $L_\infty^\ast=b-\bar x$ and $nL_\infty^\ast=nb-\sum_tx_t=L_1^\ast$. In general a clipped uniform raise satisfies the constraint with total at most $nL_\infty$, hence $L_1^\ast\le nL_\infty^\ast$. \qed

\paragraph{Proposition~\ref{prop:min} (minimum).} $\min_ty_t\ge b$ iff $y_t\ge b$ for all $t$; raising task $t$ to exactly $\max(x_t,b)$ is cheapest and possible iff $c_t\ge b$ whenever $x_t<b$; since $c_t\ge x_t$, this is the same as $c_t\ge b$ for all $t$. Tasks with $x_t\ge b$ are untouched, and the costs follow. \qed

\paragraph{Proposition~\ref{prop:median} (median, $n$ odd).} With $m=(n+1)/2$, $\mathrm{median}(y)\ge b$ iff at least $m$ coordinates are $\ge b$. A task already at $x_t\ge b$ counts at no cost; every other task can be lifted to exactly $b$ at cost $b-x_t$ if $c_t\ge b$. The cheapest way to have $m$ coordinates $\ge b$ lifts the tasks with the smallest shortfalls $(b-x_t)_+$ among those with $c_t\ge b$, which is $L_1^\ast$; $K^\ast=\max(0,m-|\{t:x_t\ge b\}|)\le m$. Feasibility needs at least $m$ tasks with $c_t\ge b$. \qed

\paragraph{Monotonicity.} Every rule here is nondecreasing in the target's own scores, so the set of feasible raises within the cap box is an up-set, and first place is reachable under cap $c$ iff the target raised to $c$ is first.

\section{Additional results}
\label{app:results}

\begin{table}[h]
\caption{Closed forms re-derived independently of the solvers against solver output on HELM (79 models, cap 100): maximum absolute difference per rule; the number in brackets is the number of targets to which the uniform-raise form of $L_\infty$ applies.}
\label{tab:theory}
\centering\small
\input{../figures/tables/tab_theory}
\end{table}

\begin{table}[h]
\caption{$R^2$ of $\log L_1$ on $\log$ rank and on $\log$ gap (ceiling cap; targets with positive gap and cost), with 95\% bootstrap intervals over targets (descriptive). The last column is the difference of the two $R^2$ with its interval; for the mean the gap fit is an identity. Ordinal rules on the Open LLM Leaderboard use the rank-stratified sample; plurality there is untestable (rank ties).}
\label{tab:c4}
\centering\scriptsize
\resizebox{\linewidth}{!}{\input{../figures/tables/tab_c4}}
\end{table}

\paragraph{Minimum versus mean by dispersion.} Table~\ref{tab:c2} gives the median cost ratio $L_1^{\min}/L_1^{\mathrm{mean}}$ in terciles of the target's coefficient of variation, with 95\% bootstrap intervals over targets, at the ceiling and under the empirical caps (Open LLM Leaderboard; the four HELM targets that are feasible under the empirical caps are omitted).

\begin{table}[h]
\caption{Median $L_1^{\min}/L_1^{\mathrm{mean}}$ by tercile of the coefficient of variation (CV) of the target's task scores; ``min cheaper'' is the share of targets for which the minimum rule needs fewer points; Spearman correlation with CV over all targets.}
\label{tab:c2}
\centering\scriptsize
\input{../figures/tables/tab_c2}
\end{table}

\begin{table}[t]
\caption{Minimum investment per rule, dataset and cap: feasible targets/targets, then median (Q1--Q3) tasks touched $K$ and points $L_1$ over \emph{feasible} targets only, so medians of rows with different feasible counts refer to different sets of targets and are not comparable. OLLB columns use the 102-target rank-stratified sample, so their medians are not population medians. Empirical 95\% caps: per-task 95th percentile of observed Open LLM Leaderboard fine-tuning gains, applied to HELM through the MMLU-Pro value. The trimmed mean drops $t=1$ of 6 scores at each end on OLLB and $t=5$ of 57 on HELM (36 targets), so its two lines are not comparable.}
\label{tab:main}
\centering\scriptsize
\resizebox{\linewidth}{!}{\input{../figures/tables/tab_main}}
\end{table}

\paragraph{Cap realism.} Table~\ref{tab:c8} gives, per rule, the share of targets for which first place is infeasible under the per-task empirical 95\% caps and the share whose $h^\ast$ exceeds the pooled 95th percentile of fine-tuning gains (\empPooled{} points). Ordinal rules on the Open LLM Leaderboard are over the 102-target rank-stratified sample in the all-target block ($\dagger$), not over all 4495 targets.

\begin{table}[h]
\caption{Share infeasible under the per-task empirical 95\% cap (95\% Wilson interval computed from the share and the number of targets; the interval is not degenerate when all targets agree, e.g.\ about [0, 0.16] for 0 of 20) and share of targets whose $h^\ast$ exceeds the pooled empirical 95th percentile; all targets, then ranks 11--30. $\dagger$: Open LLM Leaderboard ordinal rules, over the 102-target rank-stratified sample (weighted, over-representing top ranks; one Borda target unsolved), not comparable with the other rows of that block; in the ranks 11--30 block all rules use the same targets. HELM caps apply the MMLU-Pro per-task value (\empMmluQ{} points) to every subject.}
\label{tab:c8}
\centering\scriptsize
\input{../figures/tables/tab_c8}
\end{table}

\paragraph{Minimum cap by rank, task-subset scaling and robustness runs.} Figure~\ref{fig:hstarbins} shows the median $h^\ast$ by rank, Figure~\ref{fig:tasksub} the task-subset scaling of Section~\ref{sec:exp-theory} and Figure~\ref{fig:robust} the noise and task-resampling runs of Section~\ref{sec:exp-robust}.

\begin{figure}[h]
\centering
\includegraphics[width=0.95\linewidth]{../figures/fig_hstar_bins}
\caption{Median minimum uniform cap $h^\ast$ (points) by the target's rank under the mean, per rule, with 95\% bootstrap intervals over targets; targets per bin in brackets. Reference lines as in Figure~\ref{fig:cap}.}
\label{fig:hstarbins}
\end{figure}

\begin{figure}[h]
\centering
\includegraphics[width=0.8\linewidth]{../figures/fig_tasksub}
\caption{$K/n$ for random subsets of $n$ of the 57 HELM subjects (50 draws per $n$, cap $+10$), for the models ranked 2 (left) and 5 (right). Line: median over feasible draws; band: bootstrap interval of the median.}
\label{fig:tasksub}
\end{figure}

\begin{figure}[h]
\centering
\includegraphics[width=0.95\linewidth]{../figures/fig_robust}
\caption{Noise and task resampling on rank-stratified samples (cap $+10$; HELM filled, Open LLM Leaderboard open; numbers of targets and draws are in the panel titles, HELM evaluation noise has 100 draws). Bars are 95\% intervals across draws; the number of targets entering each rule's ratio (HELM$\mid$Open LLM Leaderboard; ``-'': not run) is under the rule name. (a) Ratio of $L_1$ under evaluation noise to $L_1$ at the observed scores. (b) The same ratio under a bootstrap over HELM subjects. (c) Share of target-draws that are feasible under noise (marker) against the observed scores (dash). (d) Extra raise beyond $\varepsilon$ for a just-first target to stay first in 95\% of noise draws (ranks 11--30). The Condorcet and plurality rules were not run for (d), and the ordinal rules were not run on the Open LLM Leaderboard here, hence the missing symbols.}
\label{fig:robust}
\end{figure}

\paragraph{Other caps.}
Table~\ref{tab:caps} reports the front+0, front+2, $+5$, $+20$ and empirical 99\% caps in the layout of Table~\ref{tab:main}.

\begin{table}[h]
\caption{Minimum investment under the remaining caps, same layout as Table~\ref{tab:main}. The trimmed mean was not run under these caps (n/a).}
\label{tab:caps}
\centering\scriptsize
\resizebox{\linewidth}{!}{\input{../figures/tables/tab_caps_appendix}}
\end{table}

\section{Robustness details}
\label{app:robust}

\paragraph{Margin $\varepsilon$.} A margin of $10^{-6}$ instead of $10^{-3}$ changes no feasibility and no median cost (largest per-group ratio \epsSmallMaxRatioHelm{} on HELM and \epsSmallMaxRatioOllb{} on the Open LLM Leaderboard). A margin of $0.1$ raises the median $L_1$ by factors \epsRatioLoHelm--\epsRatioHiHelm{} (HELM) and \epsRatioLoOllb--\epsRatioHiOllb{} (Open LLM Leaderboard) across rules at the ceiling and cap $+10$, and turns at most \epsFeasChangeHelm{} HELM and \epsFeasChangeOllb{} Open LLM Leaderboard targets from feasible to infeasible in any (rule, cap) cell; single targets whose base cost is small can change by much larger factors.

\paragraph{Slices of the Open LLM Leaderboard.} Table~\ref{tab:slices} compares the full leaderboard (ceiling cap, rank-stratified sample: 20 targets of rank 11--30, 70 of rank 31--100) with the merge-free slice, the raw (un-normalised) scores and the top-500 slice. Dropping merged models moves most medians by a few points at most (ranks 11--30: from $-1.7$ for the minimum to $+2.1$ for the Condorcet requirement, $+0.9$ for the mean; ranks 31--100: up to $+3.5$ for the median, $-2.8$ for the minimum) but the order of the median and the minimum rule is not robust to the slice (next paragraph). For the score rules the top-500 slice equals the full leaderboard (a target's cost depends only on the models above it), and only the ordinal rules differ, by up to 3 points. The raw scale changes the picture: the median rule's cost in ranks 11--30 rises from 4.2 to 20.8 points while the mean's falls from 32.2 to 23.0, so the normalisation matters for rules that use the order or the spread of task scores. Intervals are bootstrap intervals over the sampled targets; they do not cover the choice of slice. Cap $+10$ runs exist for all three slices (summarised in the released slice file) but are not tabulated here.

\paragraph{Median versus minimum across slices.} The slice medians use different target sets: over the sampled targets of rank 2--100 (ceiling cap; \sliceNMain{} targets in the main sample and the merge-free and raw slices, \sliceNTop{} in the top-500 slice) the median $L_1$ of the median rule is \sliceLMainMedian{} and that of the minimum rule \sliceLMainMin{} in the main analysis, so the median rule is cheaper; without merged models the two are \sliceLNomergeMedian{} and \sliceLNomergeMin{}, on raw scores \sliceLRawMedian{} and \sliceLRawMin{}, and in the top-500 slice \sliceLTopMedian{} and \sliceLTopMin{}, so the minimum is cheaper in all three. A statement that one of these two rules is cheaper on the Open LLM Leaderboard is therefore conditional on the data slice (the mean rule costs \sliceLMainMean, \sliceLNomergeMean, \sliceLRawMean{} and \sliceLTopMean{} in the four slices).

\begin{table}[h]
\caption{Open LLM Leaderboard slices, rank bin 11--30, ceiling cap: median $K$ and $L_1$ by rule.}
\label{tab:slices}
\centering\scriptsize
\resizebox{\linewidth}{!}{\input{../figures/tables/tab_slices}}
\end{table}

\paragraph{Spillover (stretch).} If effort on one task lowers the other tasks by $\kappa(D-d_t)$ with $\kappa=\lambda/(n-1)$, the exact solvers give a median $L_1$ ratio to the no-spillover cost of $\spillHelmMeanLamA$ ($\lambda=0.1$) and $\spillHelmMeanLamB$ ($\lambda=0.25$) for the mean, $\spillHelmMinLamA$ and $\spillHelmMinLamB$ for the minimum on HELM, and $\spillOllbMinLamA$, $\spillOllbMinLamB$ on the Open LLM Leaderboard; for the mean the ratio equals $1/(1-\lambda)$. Spillover therefore raises costs most for the mean and least for the minimum. The solver at $\lambda=0$ reproduces the main solvers.

\section{Further sensitivity analyses}
\label{app:sens}

These analyses repeat parts of the main grids under changed assumptions; none changes the headline comparison of the minimum with the mean, and two change the rule order under gain caps. Targets are those of the main grids (HELM: 78 non-leaders; Open LLM Leaderboard: ranks 2--100, \cvxOllbLinearN{} targets, ceiling cap, margin $\varepsilon=10^{-3}$ unless stated).

\paragraph{Convex task cost.} The main grids charge one unit per score point. We re-solved the minimum-cost problem exactly for the mean and the minimum under a quadratic cost $(y_t-x_t)^2$ (a large jump is dearer) and a log-headroom cost $\ln\frac{100-x_t}{100-y_t}$ (points near the ceiling are dearer); the linear column reproduces the main-grid costs. Median ratio $L_1^{\min}/L_1^{\mathrm{mean}}$ (linear / quadratic / log-headroom): HELM \cvxHelmLinearRatio{} / \cvxHelmQuadraticRatio{} / \cvxHelmLogheadRatio, Open LLM Leaderboard \cvxOllbLinearRatio{} / \cvxOllbQuadraticRatio{} / \cvxOllbLogheadRatio. The minimum is cheaper than the mean for \cvxHelmLinearShare\% / \cvxHelmQuadraticShare\% / \cvxHelmLogheadShare\% of HELM targets and \cvxOllbLinearShare\% / \cvxOllbQuadraticShare\% / \cvxOllbLogheadShare\% of Open LLM Leaderboard targets. The finding therefore holds for most targets but is not universal under a quadratic cost: on the Open LLM Leaderboard the mean is cheaper for \cvxOllbQuadHighMeanCheaper\% of the high-dispersion tercile (median ratio there \cvxOllbQuadHighRatio). That the ratio rises with dispersion persists (Spearman correlation with the coefficient of variation: HELM \cvxHelmLinearRho{} linear, \cvxHelmQuadraticRho{} quadratic; Open LLM Leaderboard \cvxOllbLinearRho{} and \cvxOllbQuadraticRho). Costs are stylised and equal across tasks; this tests the shape of the cost, not task-specific prices.

\paragraph{Noise-robust margins.} Because $\varepsilon=10^{-3}$ is not a robust margin (Section~\ref{sec:exp-robust}), we re-ran the comparison at margins $0.1$ and $1$ point and at rule-specific margins (2 points for the mean and the median, 8 for the minimum, about the extra margin each needs for a 95\% stay probability). The minimum remains cheaper than the mean for every target in all four settings; its median cost ratio is \mrHelmBaseRatio{} (base), \mrHelmTenthRatio{} ($0.1$), \mrHelmOneRatio{} ($1$) and \mrHelmRuleRatio{} (rule-specific) on HELM and \mrOllbBaseRatio, \mrOllbTenthRatio, \mrOllbOneRatio{} and \mrOllbRuleRatio{} on the Open LLM Leaderboard, so the advantage roughly halves with rule-specific margins, and the correlation with dispersion stays high (\mrHelmRuleRho{} and \mrOllbRuleRho). The rule order under fixed gains changes: at cap $+10$ the share of targets for which first is feasible is, under the base margin and rule-specific margins, \mrFeasHelmBaseMean\% and \mrFeasHelmRuleMean\% (mean), \mrFeasHelmBaseMedian\% and \mrFeasHelmRuleMedian\% (median), \mrFeasHelmBaseMin\% and \mrFeasHelmRuleMin\% (minimum) on HELM, and \mrFeasOllbBaseMean\% and \mrFeasOllbRuleMean\%, \mrFeasOllbBaseMedian\% and \mrFeasOllbRuleMedian\%, \mrFeasOllbBaseMin\% and \mrFeasOllbRuleMin\% on the Open LLM Leaderboard. The minimum, which needs the largest margin, is hardly ever feasible under rule-specific margins, and on the Open LLM Leaderboard the median overtakes the mean. Statements about which rule is most or least feasible under fixed gains are thus conditional on the margin.

\paragraph{Dependence between targets: clusters by base model.} The Open LLM Leaderboard targets include fine-tunes and merges of one another. We resampled whole clusters defined by the user-entered \texttt{Base Model} field (\cbClusters{} clusters for \cbTargets{} targets; models without a usable entry form their own cluster; 1000 resamples). The median cost ratio of the minimum to the mean is \cbRatio{} with a cluster interval [\cbRatioLo, \cbRatioHi] against a target-level interval [\cbRatioIidLo, \cbRatioIidHi], and \cbRatioHigh{} [\cbRatioHighLo, \cbRatioHighHi] in the high-dispersion tercile. The difference of $R^2$ of $\log L_1$ on $\log$ gap and on $\log$ rank, with cluster intervals, is \cbDiffGeomean{} [\cbDiffGeomeanLo, \cbDiffGeomeanHi] (geometric mean), \cbDiffHarmonic{} [\cbDiffHarmonicLo, \cbDiffHarmonicHi] (harmonic), \cbDiffMedian{} [\cbDiffMedianLo, \cbDiffMedianHi] (median) and \cbDiffMin{} [\cbDiffMinLo, \cbDiffMinHi] (minimum); for Borda and the Condorcet requirement (ordinal sample) the intervals include zero. Clusters widen the intervals only slightly and change no conclusion. The field is user-entered: \bmPlaceholder{} of the 540 fine-tune/base pairs of the empirical gains have a base entry that is a placeholder string or does not resolve, so some clusters are incomplete, and \bmSuspect{} pairs have the fine-tune below its base on all six tasks. HELM has no such field, so its intervals stay at the level of targets.

\paragraph{Composition of the empirical-gain pairs.} Of the 540 pairs, \pairDomain{} (\pairDomainShare\%) are fine-tuned on domain-specific data and \pairChat{} (\pairChatShare\%) are chat, RLHF or preference-tuned models, which are trained for several tasks at once and may not be informative about a developer who targets one task. Dropping chat-typed models, merges and fine-tunes whose name signals instruction tuning leaves \pairRestrictedN{} pairs; the pooled 95\% quantile falls from \empPooled{} to \empPooledRestricted{} points, and the shares of targets with first place infeasible under the per-task caps barely change (ranks 11--30 on HELM: mean \rpHelmMeanA\% to \rpHelmMeanR\%, median \rpHelmMedianA\% to \rpHelmMedianR\%, minimum \rpHelmMinA\% to \rpHelmMinR\%; ranks 2--30 on the Open LLM Leaderboard: mean \rpOllbMeanA\% to \rpOllbMeanR\%, median \rpOllbMedianA\% to \rpOllbMedianR\%, minimum \rpOllbMinA\% to \rpOllbMinR\%).

\paragraph{Ordinal rules for every target of ranks 2--100.} In the main text the Open LLM Leaderboard ordinal rules use a stratified sample; here we ran Borda, plurality and the Condorcet requirement for all \cvxOllbLinearN{} targets of ranks 2--100. At the ceiling every target is feasible for each of them; the median number of tasks touched and the median $L_1$ are Borda \ordHundredCeilBordaK{} tasks and \ordHundredCeilBordaL{} points, plurality \ordHundredCeilPluralityK{} and \ordHundredCeilPluralityL, Condorcet \ordHundredCeilCondorcetK{} and \ordHundredCeilCondorcetL, against \ordHundredCeilMeanK{} and \ordHundredCeilMeanL{} for the mean and \ordHundredCeilMinK{} and \ordHundredCeilMinL{} for the minimum. At cap $+10$ the share with first place feasible is Borda \ordHundredTenBordaFeas\%, Condorcet \ordHundredTenCondorcetFeas\%, plurality \ordHundredTenPluralityFeas\% and, for the score rules, mean \ordHundredTenMeanFeas\%, geometric mean \ordHundredTenGeomeanFeas\%, median \ordHundredTenMedianFeas\%, harmonic mean \ordHundredTenHarmonicFeas\%, minimum \ordHundredTenMinFeas\%; plurality is the hardest rule to make feasible under a fixed gain on this dataset, as in the sample.

\paragraph{Computation.} The released jobs used about \cpuHours{} CPU-hours as metered by the scheduler (summed over the jobs finished at the time of the last snapshot; the host was shared), about five times the 4 CPU-hour estimate in the task description. The largest solve for the ordinal rules peaked at \peakMemMB{} MB resident memory.

